\documentclass[10pt,letterpaper]{article} % exam class defaults to 10pt on letterpaper
\usepackage[margin=1in]{geometry}         % exam class defaults to 1-inch margins
\usepackage{graphicx} 
\usepackage{fancyhdr}
\usepackage{amssymb}
\usepackage{amsmath}
\usepackage{mathtools}
\usepackage[most]{tcolorbox}
\usepackage[dvipsnames]{xcolor}
\usepackage{blindtext}
\usepackage{amsthm}
\usepackage{upgreek}
\usepackage{hyperref}
\usepackage{tikz}
\usetikzlibrary{positioning}
%making an edit, want to see if auto script commits it%
%making another edit to test the commit bot

\setlength{\parindent}{0pt}
\setlength{\parskip}{1ex plus 0.5ex minus 0.2ex} % Matches the paragraph spacing of exam
\hbadness=10000 % Place in your preamble


\newtheorem{thm}{Theorem}[section]
\newtheorem{lem}{Lemma}[section]
\newtheorem{defn}{Definition}[section]
\newtheorem{cor}{Corollary}[section]
\newtheorem{axm}{Axiom}[section]

\tcolorboxenvironment{thm}{colback=lime, colframe=black, breakable}
\tcolorboxenvironment{lem}{colback=yellow, colframe=black, breakable}
\tcolorboxenvironment{defn}{colback=pink, colframe=black, breakable}
\tcolorboxenvironment{cor}{colback=SpringGreen, colframe=black, breakable}
\tcolorboxenvironment{axm}{colback=white, colframe=black, breakable}

% --- Header/Footer Setup ---
\fancyhf{}
\renewcommand\headrulewidth{0.5pt}
\pagestyle{fancy}
\author{Omkar Tasgaonkar}
\setlength{\headsep}{0.5in}
\fancyfoot[C]{(\thepage)}
\renewcommand\footrulewidth{0.5pt}

\title{Robotics}
\author{Omkar Tasgaonkar \\ \small otasgaonkar@ucla.edu}
\date{}

\fancypagestyle{firstpage}{
    \fancyhf{}
    \fancyhead[C]{\textbf{Notes in Engineering}}
    \renewcommand\headrulewidth{0.5pt}
    \fancyfoot[C]{(\thepage)}
    \renewcommand\footrulewidth{0.5pt}
}
\begin{document}
\maketitle
This is a set of notes covering interesting concepts I find while studying the Modern Robotics textbook. I will try to bring analytic theory whenever applicable.

\begin{tcolorbox}[colback=white, colframe=black, arc=0mm]
Textbooks Used to Develop these Notes:
\begin{enumerate}
\item \textit{\textbf{Modern Robotics} - Kevin Lynch, Frank Park}
\item \textit{\textbf{Calculus on Manifolds} - Michael Spivak}
\item \textit{\textbf{Mathematical Analysis} - Tom Apostol}
\end{enumerate}
\end{tcolorbox}
\begin{tcolorbox}[colback=white, colframe=black, arc=0mm]
\textbf{Notes Last Updated:} \today
\end{tcolorbox}
\thispagestyle{firstpage}
\newpage
\begin{tcolorbox}[colback=white, colframe=black, arc=0mm]
\tableofcontents
\end{tcolorbox}
\textbf{NOTE:} Any sections/subsections marked with an asterisk are extra sections and have been added to this document as extensions of introduced theory.
\pagebreak

\section{Holonomic Constraints}
\subsection{Basic Multivariable Analysis}
\textbf{Note:} The following notes may not be entirely rigorous! I have attempted to write them logically to the best of my abilities so that they follow a structured narrative.\\

We begin with the exploration of multivariate analysis. When we deal with a multivariate function in $\mathbb{R}$, we deal with some map of the form $f: \mathbb{R}^n \xrightarrow{} \mathbb{R}^m$. We can define this more formally via vector operations as the following:
\begin{equation}
f \left(\begin{bmatrix}
x_1\\.\\.\\.\\x_n
\end{bmatrix}\right) = \begin{bmatrix}
f_1(x_1, x_2, ..., x_n)\\
.\\.\\.\\
f_m(x_1, x_2, ..., x_n)
\end{bmatrix}
\end{equation}
such that $\forall j \in [1, m]: f_j: \mathbb{R}^n \xrightarrow{} \mathbb{R}$. Our original definition of the derivative in $\mathbb{R}$ was:
\begin{equation*}
f'(x) \coloneq \lim_{z \xrightarrow{} x} \frac{f(z) - f(x)}{z-x}
\end{equation*}
or equivalently
\begin{equation*}
\lim_{z \xrightarrow{} x} \frac{f(z) - f(x) - f'(x)(z-x)}{z-x} =0
\end{equation*}
Using the substitution $h = z-x$, we also get the equivalent definition:
\begin{equation*}
\lim_{h \xrightarrow{} 0}\frac{f(x+h) - f(x) - f'(x)h}{h} = 0
\end{equation*}
In which case, our $\epsilon-\delta$ definition would be:
\begin{equation*}
\forall \epsilon > 0, \exists \delta > 0, \forall h \in B(0, \delta): \frac{|f(x+h) - f(x) - f'(x)h|}{|h|} < \epsilon
\end{equation*}
However if we try to plainly substitute the definition for a function $f: \mathbb{R}^n \xrightarrow{} \mathbb{R}^m$ into the first definition, we get "division of two vectors" which is not a defined operation for vector spaces (they are only abelian groups under addition).\\

We consider the following lemma in order to understand how to extend the derivative to the multivariate case:
\begin{lem}
$f'(x)$ exists $\iff \left\{\frac{1}{\delta}\sup_{z \in B(x, \delta)\setminus \{x\}}|f(z) - f(x) - f'(x)(z-x)|\right\}_{\delta > 0} \xrightarrow{} 0$\\

\end{lem}
\begin{proof}
$\implies:$ If $f'(x)$ exists, then:
\begin{equation*}
\forall \epsilon > 0, \exists \delta_0 > 0, \forall z \in B(x, \delta_0) \setminus \{x\}: \left|\frac{f(z) - f(x)}{z-x} - f'(x)\right| < \epsilon
\end{equation*}
\begin{equation*}
\implies \forall \epsilon >0, \exists \delta_0 > 0, \forall \delta \leqslant \delta_0, \forall z \in B(x, \delta)\setminus \{x\}: \left|\frac{f(z) - f(x)}{z-x} - f'(x)\right| < \epsilon
\end{equation*}
\begin{equation*}
\implies |f(z) -f(x) - f'(x)(z-x)| < \epsilon \cdot |z-x| < \epsilon \cdot \delta
\end{equation*}
Because this is true for all $z$ in the punctured $\delta$-ball around $x$, we have an upper bound, so we take the supremum:
\begin{equation*}
\implies \sup_{z \in B(x, \delta) \setminus \{x\}}|f(z) - f(x) - f'(x)(z-x)| < \epsilon \cdot \delta
\end{equation*}
And so we rearrange the terms to get a $\delta$-net that converges such that:
\begin{equation*}
\forall \epsilon > 0, \exists \delta_0 > 0, \forall \delta \leqslant \delta_0: \frac{1}{\delta}\sup_{z \in B(x, \delta)\setminus \{x\}}|f(z) - f(x) - f'(x)(z-x)| < \epsilon
\end{equation*}
Thus proving the forward direction.\\

$\impliedby$: Now for the converse, consider:
\begin{equation*}
\forall \epsilon >0, \exists \delta_0 > 0, \forall \delta \leqslant \delta_0: \sup_{z \in B(x, \delta)\setminus\{x\}}|f(z) - f(x) - f'(x)(z-x)| < \epsilon \cdot \delta
\end{equation*}
Now for any arbitrary $\delta$-ball around $x$ where $\delta \leqslant \delta_0$, we have that:
\begin{equation*}
\sup_{z \in B(x, \delta)\setminus \{x\}}|f(z) - f(x) - f'(x)(z-x)| < \epsilon \cdot \delta
\end{equation*}
\begin{equation*}
\implies \forall z \in B(x, \delta) \setminus \{x\}: \frac{1}{\delta}|f(z) - f(x) - f'(x)(z-x)| < \epsilon
\end{equation*}
\begin{equation*}
\implies \forall z \in B(x, \delta_0) \setminus B(x, \delta): \frac{|f(z) - f(x) - f'(x)(z-x)|}{|z-x|} < \frac{1}{\delta}|f(z) - f(x) - f'(x)(z-x)| < \epsilon
\end{equation*}
Since $\delta$ was arbitrary as long as it was less than or equal to $\delta_0$, we observe:
\begin{equation*}
\forall \epsilon > 0, \exists \delta_0 > 0, \forall \delta \leqslant \delta_0, \forall z \in B(x, \delta_0)\setminus B(x, \delta): \left|\frac{f(z) - f(x)}{z-x} - f'(x)\right| < \epsilon
\end{equation*}
\begin{equation*}
\iff \forall \epsilon > 0, \exists \delta_0 > 0, \forall z \in \bigcup_{\delta \in (0, \delta_0]}(B(x, \delta_0)\setminus B(x, \delta)): \left|\frac{f(z) - f(x)}{z-x} - f'(x)\right| < \epsilon
\end{equation*}
\begin{equation*}
\iff \forall \epsilon > 0, \exists \delta_0 > 0, \forall z \in B(x, \delta_0) \setminus \bigcap_{\delta \in (0, \delta_0]}B(x, \delta): \left|\frac{f(z) - f(x)}{z-x} - f'(x)\right| < \epsilon
\end{equation*}
Now because $\mathbb{R}$ is a $T_1$ space, we see that the intersection in the previous statement is equivalent to $\{x\}$, so we get:
\begin{equation*}
\forall \epsilon > 0, \exists \delta_0 > 0, \forall z \in B(x, \delta_0) \setminus \{x\}: \left|\frac{f(z) - f(x)}{z-x} - f'(x)\right| < \epsilon
\end{equation*}
A useful equivalent statement is:
\begin{equation*}
\forall \epsilon > 0, \exists \delta_0 < \epsilon, \forall \delta \leqslant \delta_0, \forall z \in B(x, \delta) \setminus \{x\}: \frac{1}{\delta}|f(z) - f(x) - f'(x)(z-x)| < \epsilon
\end{equation*}
\end{proof}
Now essentially this supremum $\delta$-net, similar to oscillation is the worst case net. The important part of this result is the idea of linear approximation, such that:
\begin{equation*}
f'(x)(z-x) + f(x) -\epsilon^2 < f(z) < f'(x)(z-x) + f(x) + \epsilon^2
\end{equation*}
or as we write it in calculus $f(z) \approx f(x) + f'(x)(z-x)$, which is the tangent line approximation.\\

We can apply this approximation idea to the multivariate case, but first we introduce the following notion:

\medskip
\begin{defn}
Given some $f: \mathbb{R}^n \xrightarrow{} \mathbb{R}^m$ in the form of equation $(1)$, then for some $i \in [1, m], j \in [1, n]$:
\begin{equation*}
\frac{\partial f_i(x)}{\partial e_j} \coloneq \lim_{h \xrightarrow{} 0}\frac{f_i(x + e_jh) - f_i(x)}{h}
\end{equation*}
where $\{e_1,\cdots, e_n\}$ is the canonical basis of $\mathbb{R}^n$. This is the definition of the partial derivative.
\end{defn}
Since each $f_i: \mathbb{R}^n \xrightarrow{} \mathbb{R}$, we see that the partial derivative is scalar valued. We now show that it admits the linear approximation property. We can rewrite our limit definition as:
\begin{equation*}
\lim_{h \xrightarrow{} 0} \left(\frac{f_i(x+e_jh) - f_i(x)}{h} - \frac{\partial f_i(x)}{\partial e_j}\right) = \lim_{h \xrightarrow{} 0} \frac{f_i(x+e_jh) - f_i(x) - (\frac{\partial f_j(x)}{\partial e_j})h}{h} = 0
\end{equation*}
So with an argument similar to Lemma 1.1, we see that $f_i(z) \approx f_i(x) + \left(\frac{\partial f_i}{\partial e_j}\right)h$. Now observe that:
\begin{equation*}
f(z) \coloneq \begin{bmatrix}
f_1(z)\\.\\.\\.\\f_m(z)
\end{bmatrix}
\end{equation*}
\begin{lem}
For all $i \in [1, m]$ and some $x \in \mathbb{R}^n$ and an $e_j$:
\begin{equation*}
\frac{\partial f_i(x)}{\partial e_j} \ \text{exist}
\end{equation*}
if and only if:
\begin{equation*}
\frac{\partial f(x)}{\partial e_j} \ \text{exists}
\end{equation*}
\end{lem}
\begin{proof}
$\implies:$ We claim that the partial derivative of $f$ takes the following form:
\begin{equation*}
 \frac{\partial f(x)}{\partial e_j} = \begin{bmatrix}
\frac{\partial f_1(x)}{\partial e_j}\\.\\.\\.\\\frac{\partial f_m(x)}{\partial e_j}\end{bmatrix}
\end{equation*}
We have some $\delta_i$ given by the existence of each partial derivative, so we define $\delta_0 \coloneq \min\{\delta_i : 1 \leqslant i \leqslant n\}$.
Then for all $h \in B(0, \delta_0)$:
\begin{equation*}
\frac{||f(x+e_jh) - f(x) - (\frac{\partial f(x)}{\partial e_j})h||}{|h|} \leqslant \frac{1}{|h|}\sum_{i=1}^{m}\left|f_i(x+e_jh) - f_i(x) - \left(\frac{\partial f_i(x)}{\partial e_j}\right)h\right| < m\epsilon
\end{equation*}
where we used the inequality that the $||v|| \leqslant \sum |v_i|$.\\

$\impliedby:$ We claim that:
\begin{equation*}
\frac{\partial f_i(x)}{\partial e_j} = e_i \cdot \frac{\partial f(x)}{\partial e_j}
\end{equation*}
We now use the inequality $|v_i| \leqslant ||v||$ and find that $\forall i \in [1, m]$:
\begin{equation*}
\frac{|f_i(x+e_jh) - f_i(x) - (e_i \cdot \frac{\partial f(x)}{\partial e_j})h|}{|h|} = \frac{|e_i \cdot f(x+e_jh) - e_i \cdot f_i(x) - (e_i \cdot \frac{\partial f(x)}{\partial e_j})h|}{|h|}
\end{equation*}
\begin{equation*}
= \frac{|e_i \cdot (f(x+e_jh) - f(x) - \frac{\partial f(x)}{\partial e_j}h)|}{|h|}\leqslant \frac{||f(x+e_jh) - f(x) - (\frac{\partial f(x)}{\partial e_j})h||}{|h|} < \epsilon
\end{equation*}
so the same $\delta$ suffices.\\

Now once again with an argument similar to Lemma 1.1, we can extend the idea of linear approximation and say that for all $z$ such that $z = x + e_jh$:
\begin{equation*}
f(z) \approx f(x) + \left(\frac{\partial f(x)}{\partial e_j}\right)h
\end{equation*}
\end{proof}
Notice that the proofs did not use any properties of the basis vector $e_j$ itself, so we can extend the definition:
\begin{defn}
Given some $f: \mathbb{R}^n \xrightarrow{} \mathbb{R}^m$ in the form of equation $(1)$, then for some $v \in \mathbb{R}^n$ and $i \in [1, m], j \in [1, n]$:
\begin{equation*}
\frac{\partial f_i(x)}{\partial v} \coloneq \lim_{h \xrightarrow{} 0}\frac{f_i(x+hv) - f_i(x)}{h}
\end{equation*}
and similarly for $f$. Such a derivative is called the Gateaux derivative.
\end{defn}
We can extend the Gateaux derivative to obtain the Frechet derivative. The Gateaux derivative limited itself to a specific direction of the vector $v$ (and then used $h$ as the variable for the limit).
\begin{defn}
The Frechet derivative of a function $f$ at $x \in \mathbb{R}^n$ is:
\begin{equation*}
Df(x): \lim_{v \xrightarrow{} 0} \frac{||f(x + v) - f(x) - (Df(x))(v)||}{||v||} = 0
\end{equation*}
or equivalently:
\begin{equation*}
\forall \epsilon > 0, \exists \delta > 0, \forall v \in B(0, \delta): \frac{||f(x+v) - f(x) - (Df(x))(v)||}{||v||} < \epsilon
\end{equation*}
where the norm induces a metric so that $v \in B(0, \delta)$ (in this case $0$ is the $0$ vector), is the same as saying $||v|| \in B(0, \delta)$ or equivalently $||v|| < \delta$.
\end{defn}
The form of the limit is similar to the idea of linear approximation. Since $v \in \mathbb{R}^n$, then we must have that $Df(x) \in \mathcal{M}_{m \times n}(\mathbb{R})$. If the function is Frechet differentiable, then $Df(x)$ is the Jacobian:
\begin{equation*}
Df(x) = \begin{bmatrix}
\frac{\partial f_1(x)}{\partial e_1} & \cdots &  \frac{\partial f_1(x)}{\partial e_n}\\
\vdots & \ddots & \vdots\\
\frac{\partial f_m(x)}{\partial e_1} & \cdots & \frac{\partial f_m(x)}{\partial e_n} 
\end{bmatrix} = \begin{bmatrix}
\frac{\partial f(x)}{\partial e_1} & \cdots & \frac{\partial f(x)}{\partial e_n}
\end{bmatrix}
\end{equation*}
\begin{proof}
If we chose the family of vectors $\{he_j : h \in \mathbb{R} \land j \in [1, n]\}$, then differentiability tells us:
\begin{equation*}
\forall \epsilon > 0, \exists \delta > 0, \forall h \in B(0, \delta): \frac{||f(x+he_j) - f(x) - h(Df(x))e_j||}{|h|} < \epsilon
\end{equation*}
telling us that:
\begin{equation*}
\frac{\partial f(x)}{\partial e_j} = Df(x)e_j
\end{equation*}
which is the $j$-th column of the Jacobian matrix.
\end{proof}
Notice that Lemma 1.2 was a statement on equivalence; however, the converse is not necessarily true for the Jacobian matrix. The total differentiability of $f$ at $x$ gives us the existence of partial derivatives, but not necessarily the converse. In order to prove this, we require the following concept:

%Need to expound on counterexamples more%
% Consider the function:
% \begin{equation*}
% f(x_1, x_2) = \frac{x_1x_2 + x_2}{(x_1+x_2)^4}
% \end{equation*}
% when $x \neq 0$ and $0$ when $x = 0$. Then considering the partial derivative in the $e_1$ direction, we have:
% \begin{equation*}
% \frac{}{h}
% \end{equation*}
%Need to expound on counterexamples more%
\subsubsection{MVT in the Multivariable Case}
\begin{thm}
Given a function $f: \mathbb{R} \xrightarrow{} \mathbb{R}$ that is continuous on the closed interval $[a, b]$ and differentiable on $(a, b)$, then:
\begin{equation*}
\exists z \in (a, b): f'(z) = \frac{f(b) - f(a)}{b-a}
\end{equation*}
\end{thm}
\begin{proof}
Since $f$ is continuous on a compact interval, it achieves the extrema by EVT. Now consider the function:
\begin{equation*}
h(x) = \frac{f(b) - f(a)}{b-a}(x-a) - f(x) + f(a)
\end{equation*}
which at both $x = a$ and $x = b$, we have that $h(x) = 0$, so by Rolle's Theorem, there exists some extrema within the interval, so that:
\begin{equation*}
\exists z \in (a, b): h'(z) = 0 \implies f'(z) = \frac{f(b) - f(a)}{b-a}
\end{equation*}
\end{proof}
If we try to apply this proof to a function $f: \mathbb{R}^n \xrightarrow{} \mathbb{R}$, then in order to apply Theorem 1.1, we require some compact interval given $f$ is continuous.\\

% \begin{defn}
% A topological space $X$ is disconnected if it admits a separation such that:
% \begin{equation*}
% \exists A, B \subseteq X: A, B \ \text{open} \land A \cap B = \varnothing \land A \cup B = X
% \end{equation*}
% \end{defn}
% This tells us that $A, B$ are both open and closed; so a connected set is one which has only the subsets $X, \varnothing$ as the open AND closed subsets. We require both to be open (or equivalently closed), otherwise every space would be trivially disconnected. We also require both to be nonempty.\\

% There is an additional characterization of connectedness for the subspace topology. 
Now, the following lemma will be useful in exploring MVT:
% \begin{lem}
% An open, connected subset $U \subseteq \mathbb{R}^n$ is path connected.
% \end{lem}
% \begin{proof}
% We first need to prove any open ball $B(x, r)$ is convex. Given some $\alpha \in [0, 1]$, we want that for any two points $y, z$:
% \begin{equation*}
% \alpha y + (1-\alpha)z \in B(x, r)
% \end{equation*}
% \begin{equation*}
% |\alpha y + (1-\alpha)z - x| = |\alpha(y - x) + (1-\alpha)(z-x)| \leqslant \alpha r + (1-\alpha)r = r
% \end{equation*}
% Now, consider two points $a, b$ in $U$. \textbf{This proof is a bit beyond me right now, I will revisit it at a later time.}
% \end{proof}

\begin{lem}
Any open, convex subset $U \subseteq \mathbb{R}^n$ is path connected.
\end{lem}
\begin{proof}
Openness does not play an important role here. We can simply take the continuous function:
\begin{equation*}
f(\alpha) = \alpha x + (1-\alpha)y
\end{equation*}
for some fixed $x, y \in U$. By convexity, we know $\text{Ran}(f) \subseteq U$.
\end{proof}
\begin{lem}
An open ball $B(x, r)$ is convex.
\end{lem}
\begin{proof}
Take $y, z \in B(x, r)$, then:
\begin{equation*}
|\alpha y + (1-\alpha)z - x| = |\alpha(y-x) + (1-\alpha)(z-x)| \leqslant \alpha r + (1-\alpha) r = r
\end{equation*}
\end{proof}

\begin{thm}
Given $f_i: \mathbb{R}^n \xrightarrow{} \mathbb{R}$, we have that for some open, convex subset $U \subseteq \mathbb{R}^n$ (on which $f_i$ is differentiable) and given two points $a, b$:
\begin{equation*}
\exists z \in U: Df_i(z)(b-a) = f_i(b) - f_i(a)
\end{equation*}
where $Df_i(z)$ is the gradient vector at $z$.
\end{thm}
\begin{proof}
Since $U$ is convex, we can take the function for $t \in [0, 1]$:
\begin{equation*}
h(t) = t b + (1-t) a
\end{equation*}
such that $h(0) = a$ and $h(1) = b$. Then applying Theorem 1.1 to the function $f_i \circ h$ (which too is continuous on $[0, 1]$ and differentiable on $(0, 1)$), we see that:
\begin{equation*}
\exists c \in (0, 1): (f_i \circ h)'(c) = \frac{f_i(h(1)) - f_i(h(0))}{1-0} = f_i(b) - f_i(a)
\end{equation*}
and $(f_i \circ h)'(c) = Df_i(c) \cdot h'(c) = Df_i(c)(b-a)$, where:
\begin{equation*}
Df_i(x) = \begin{bmatrix}
\frac{\partial f_i(x)}{\partial e_1} & \cdots & \frac{\partial f_i(x)}{\partial e_n}
\end{bmatrix}
\end{equation*}
We used the multivariate chain rule in this proof, where $(f \circ h)'(t) = Df(h(t))Dh(t)$. In this case, we defined:
\begin{equation*}
h(t) = \begin{bmatrix}
t b_1 + (1-t)a_1 & \cdots & t b_n + (1-t) a_n
\end{bmatrix}^T
\end{equation*}
\end{proof}

\begin{cor}
If $\exists k \in \mathbb{R}, \exists j \in [1, n]: b-a = ke_j$ and the partial derivative of $f_i$ exists in the $e_j$ direction, then:
\begin{equation*}
\exists z \in \text{span}(e_j): \frac{\partial f_i(z)}{\partial e_j}(b-a) = f_i(b) - f_i(a)
\end{equation*}
\end{cor}
\begin{proof}
Notice this corollarly does not require full differentiability of $f_i$ (which means it does not require the gradient to be the total derivative of $f$ at said point). 
\end{proof}

\begin{lem}
The function $f: \mathbb{R}^n \xrightarrow{} \mathbb{R}^m$ is totally differentiable at a point if:
\begin{enumerate}
\item The partial derivatives exist
\item The partial derivatives are continuous at that point
\end{enumerate}
\end{lem}
\begin{proof}
We want to show:
\begin{equation*}
\forall \epsilon > 0, \exists \delta > 0, \forall v \in B(0, \delta): \frac{||f(x+v) - f(x) - (Df(x))v||}{||v||} < \epsilon
\end{equation*}
The expression $(Df(x))(v)$ gives us the following vector:
\begin{equation*}
(Df(x))(v) = \begin{bmatrix}
\sum_{i=1}^{n}\frac{\partial f_1(x)}{\partial e_i}v_i\\
\vdots\\
\sum_{i=1}^{n}\frac{\partial f_m(x)}{\partial e_i}v_i
\end{bmatrix}
\end{equation*}
and we see that:
\begin{equation*}
\frac{||f(x+v) - f(x) - (Df(x))(v)||}{||v||} \leqslant \frac{1}{||v||}\sum_{j=1}^{m}\left|f_j(x+v) - f_j(x) - \sum_{i=1}^{n}\frac{\partial f_j(x)}{\partial e_i}v_i\right|
\end{equation*}
Continuity of the partial derivatives give us:
\begin{equation*}
\forall \epsilon > 0, \exists \delta_{ji} > 0, \forall v \in B(0, \delta_{ji}): \left|\frac{\partial f_j(x+v)}{\partial e_i} - \frac{\partial f_j(x)}{\partial e_i}\right| < \epsilon
\end{equation*}
We set $\delta \coloneq \min\{\delta_{ji} : 1 \leqslant j \leqslant m \land 1 \leqslant i \leqslant n\}$. Then given some $v \in B(0, \delta)$, by Lemma 1.4, we have that the open ball is convex, so we can apply MVT to each pair $f_j(x+v) - f_j(x)$.\\

Since
\end{proof}
\subsubsection{Multivariable Chain Rule}
\subsection{Inverse and Implicit Function Theorems}
\subsection{Configuration Spaces}
The idea of configuration spaces is intertwined with the idea of manifolds, which I have not studied, so for this section we will limit ourselves to Euclidean configuration spaces. 
\end{document}